\chapter{Estrategia y Arquitectura de Microservicios}
\label{ch:arquitectura}

\section{Análisis del sistema actual}

Según el caso de estudio, BookPoint Chile nació como una librería independiente en
Concepción y, debido al crecimiento de sus ventas, abrió nuevas sucursales en Temuco y
La Serena. Actualmente opera con un \textbf{sistema monolítico} que concentra en un
mismo aplicativo las ventas, el inventario, el despacho y la atención de clientes.

El propio caso identifica tres síntomas concretos de que ese monolito ya no escala con
el crecimiento de la operación:

\begin{itemize}
    \item \textbf{Lentitud en los procesos}, provocada por el aumento sostenido de
    pedidos tanto presenciales como por canal web.
    \item \textbf{Stock desactualizado} entre sucursales, al no existir una fuente de
    verdad de inventario que se sincronice en tiempo real entre Concepción, Temuco y
    La Serena.
    \item \textbf{Demoras en la gestión de despachos}, al no contar con un módulo
    dedicado de logística que pueda evolucionar independientemente del resto del
    sistema.
\end{itemize}

Estos síntomas son característicos de una arquitectura monolítica sometida a
crecimiento: al compartir un único proceso, una única base de datos y un único ciclo de
despliegue, cualquier cambio en el módulo de inventario obliga a redesplegar
\textit{todo} el sistema -- incluida la caja de venta presencial, que no puede
detenerse sin afectar directamente la operación de las tres tiendas físicas. Del mismo
modo, un pico de tráfico en el canal web (por ejemplo, durante una campaña de
descuentos) satura recursos que también necesita la venta presencial, porque ambos
comparten el mismo proceso.

\section{Estrategia de microservicios elegida}

Para resolver el caso propuesto se optó por una arquitectura de
\textbf{microservicios independientes}, organizados por dominio de negocio (no por
capa técnica), con los siguientes principios:

\begin{itemize}
    \item \textbf{Base de datos por servicio.} Cada microservicio posee su propia base
    de datos MySQL (Tabla~\ref{tab:bd-servicio}), eliminando el acoplamiento de datos
    que hoy sufre el sistema monolítico. El módulo de inventario puede evolucionar su
    modelo de stock sin coordinar una migración con el módulo de reseñas.
    \item \textbf{Comunicación síncrona vía REST/Feign} para las relaciones entre
    dominios que sí necesitan datos de otro servicio en el momento (por ejemplo, un
    pedido necesita confirmar que el libro existe y que hay stock antes de
    concretarse).
    \item \textbf{Descubrimiento de servicios con Eureka}, para que ningún
    microservicio necesite conocer de antemano la dirección de red de los demás --
    relevante en un escenario con tres sucursales que podrían eventualmente desplegarse
    en ubicaciones distintas.
    \item \textbf{API Gateway como fachada única}, que permite exponer una sola puerta
    de entrada hacia los distintos canales (presencial, web) sin que el cliente
    necesite conocer la topología interna de microservicios.
\end{itemize}

Esta estrategia ataca directamente los tres síntomas identificados: la lentitud se
resuelve porque cada dominio escala y se despliega de forma independiente (el catálogo,
que se consulta constantemente, no compite por recursos con el módulo de pagos); el
stock desactualizado se resuelve centralizando el inventario en un único microservicio
consultado por todos los canales y sucursales; y las demoras de despacho se resuelven
con un microservicio de envíos dedicado, capaz de evolucionar su lógica de logística sin
tocar el resto del sistema.

\section{Diseño de los microservicios}

El caso de estudio describe seis perfiles de usuario (Administrador del Sistema, Jefe
de Sucursal, Asistente de Ventas, Área de Logística, Encargado de Bodega y Clientes vía
web) con un total de más de treinta acciones posibles. La Tabla~\ref{tab:mapeo-perfiles}
mapea los diez microservicios efectivamente construidos a los perfiles del caso que
cubren.

\begin{table}[h]
\centering
\small
\begin{tabularx}{\textwidth}{lX}
\toprule
\textbf{Microservicio} & \textbf{Perfiles/acciones del caso que cubre} \\
\midrule
Catálogo & Cliente: navegar catálogo de libros. \\
Inventario & Encargado de Bodega: administrar stock central y por sucursal.
Administrador: consultar stock disponible en otras sucursales. \\
Sucursales & Jefe de Sucursal: administrar datos de la sucursal. \\
Clientes & Cliente: crear cuenta, administrar perfil y direcciones. \\
Carrito & Cliente: agregar productos al carrito. \\
Pedidos & Cliente: realizar compras online. Asistente de Ventas: registrar ventas
(equivalente digital de la venta en caja). Jefe de Sucursal: generar reportes de
ventas por sucursal y por autor. \\
Pagos & Cliente: pagar una compra online. \\
Envíos & Área de Logística: gestionar envíos de pedidos a domicilio, actualizar
estados de pedidos. \\
Notificaciones & Cliente: recibir confirmaciones y alertas asociadas a su pedido. \\
Reseñas & Cliente: dejar reseñas y calificaciones de productos. \\
\bottomrule
\end{tabularx}
\caption{Mapeo de los microservicios construidos a los perfiles del caso BookPoint Chile}
\label{tab:mapeo-perfiles}
\end{table}

\subsection{Alcance y exclusiones deliberadas}

No todas las acciones descritas en el caso están cubiertas por la implementación
actual. Se priorizó el \textbf{flujo comercial completo} (explorar catálogo, agregar al
carrito, comprar, pagar, despachar, reseñar) por sobre funciones administrativas de
soporte, que quedaron fuera de alcance de forma consciente:

\begin{itemize}
    \item Gestión de usuarios y permisos por rol para el perfil Administrador del
    Sistema (autenticación básica de clientes vía JWT sí se implementó, ver
    sección~\ref{ch:etica}).
    \item Reportes de ventas por categoría (Jefe de Sucursal) -- los reportes por
    sucursal y por autor sí se implementaron (sección~\ref{ch:crud}), pero por
    categoría exigiría antes agregar un campo de categoría/género al catálogo, que
    no existe hoy.
    \item Promociones, descuentos, convenios estudiantiles y cupones (Asistente de
    Ventas, Cliente).
    \item Emisión de boletas o facturas electrónicas y gestión de devoluciones
    (Asistente de Ventas).
    \item Transferencias de stock entre sucursales, gestión de proveedores editoriales
    y optimización de rutas de distribución (Área de Logística, Encargado de Bodega).
\end{itemize}

Esta decisión de alcance no es un descuido: dado el tiempo disponible, se priorizó
profundizar el flujo transaccional principal -- con validaciones cruzadas reales entre
servicios, manejo de errores consistente y pruebas automatizadas -- por sobre cubrir
superficialmente la totalidad de acciones descritas en el caso.

\section{Plan de migración}

Para migrar el sistema monolítico descrito en el caso hacia esta arquitectura, se
plantea un enfoque incremental de tipo \textbf{strangler fig} (extracción progresiva),
en lugar de una reescritura total (\textit{big-bang rewrite}), que sería
significativamente más riesgosa para una empresa que necesita mantener sus tres
sucursales operativas durante la transición:

\begin{enumerate}
    \item \textbf{Extraer primero los dominios de solo lectura y bajo acoplamiento
    transaccional}: catálogo e inventario. Son los más consultados (cada venta,
    presencial o web, necesita leerlos) y los que menos dependen de otros módulos para
    funcionar.
    \item \textbf{Introducir el API Gateway como fachada de transición}, de modo que
    el resto del sistema (aún monolítico en esta etapa) y los microservicios ya
    extraídos convivan detrás de una única puerta de entrada, sin que el cliente note
    el cambio.
    \item \textbf{Extraer los dominios transaccionales}: clientes, carrito, pedidos y
    pagos, en ese orden, ya que cada uno depende de que el anterior esté disponible
    como servicio (el carrito necesita clientes y catálogo; los pedidos necesitan
    catálogo e inventario).
    \item \textbf{Extraer los dominios de soporte}: sucursales, envíos, notificaciones
    y reseñas, que dependen de que pedidos ya exista como microservicio independiente.
    \item \textbf{Retirar el monolito} una vez que todos los dominios estén servidos
    por microservicios y validados en producción.
\end{enumerate}

Este orden minimiza el riesgo operativo: en cada paso, el sistema sigue siendo
funcional para las tres sucursales, y solo se apaga definitivamente el monolito cuando
ya no queda ningún flujo de negocio dependiendo de él.
